% =============================================================
%  Report: free multi-domain diffusion modeling and thermal history
%  of Win-06-01b goethite (water-saturated step heat HT-Win-06-01b).
%  Figures:  ../Figures/fig_winsor_mdd_free.py  (generated/mdd_free/absolute/)
%            ../Figures/fig_winsor_mdd_free_chapter.py  (generated/mdd_free/absolute/chapter/)
%  Values:   mdd_free_report_values.tex (auto-generated by the same script)
%  Build:    latexmk -pdf winsor_mdd_free_report.tex
% =============================================================
\documentclass[11pt]{article}
\usepackage{lmodern}
\usepackage[T1]{fontenc}
\usepackage[utf8]{inputenc}
\usepackage{graphicx,amsmath,amssymb,booktabs,array}
\usepackage[margin=1in]{geometry}
\usepackage[round]{natbib}
\usepackage{caption}
\usepackage{pdflscape}
\captionsetup{font=small,labelfont=bf}
\usepackage[colorlinks=true,linkcolor=blue,citecolor=blue,urlcolor=blue]{hyperref}
\graphicspath{{../Figures/generated/mdd_free/absolute/chapter/}{../Figures/generated/mdd_free/absolute/}}
% =============================================================
%  chapter2_derived_values.tex -- AUTO-GENERATED, DO NOT EDIT BY HAND.
%  Written by Figures/write_chapter2_macros.py from
%  Figures/generated/data/stage1_regression_cache.json (itself written by
%  Figures/prepare_regression_cache.py from the master workbook).
%
%  Regenerate after any data/regression change:
%    python3 Figures/prepare_regression_cache.py
%    python3 Figures/write_chapter2_macros.py
%  (run_all_figures.py does both automatically.)
%
%  \input{} by BOTH chapter2goethite.tex (standalone) and
%  Chapter2/chapter2.tex (dissertation chapter) so the Abstract,
%  Results, Discussion, and Conclusions all read the same numbers.
%
%  Stage 1 reference line (3He, "treatment as uncertainty" central/M
%  bound, fractional 3He yield >= 0.07, n=12 points, R2=0.9890):
%    ChTwoEa/ChTwoEaErr : activation energy (kJ/mol)
%    ChTwoTc/ChTwoTcErr : Dodson (1973) closure temperature (degC),
%      intercept-only regression uncertainty (D0/a^2 term only, Ea
%      held fixed) -- see thermochr_utils.dodson_closure_temperature_stderr_interceptonly.
%      Joint slope+intercept uncertainty is +/- 14.9 degC if ever needed.
%    ChTwoRtwo/ChTwoN   : regression R^2 and number of points
%
%  Isothermal 3He permeation Arrhenius regression (from
%  generated/data/naive_lambda_arrhenius_cache.json, n=10 sequences, R2=0.9362):
%    ChTwoPermEa/ChTwoPermEaErr : activation energy, 2 s.e. (cal/mol)
%    ChTwoPermRtwo/ChTwoPermN   : regression R^2 and number of sequences
%    ChTwoPermLamMin/ChTwoPermLamMax : range of the 10 per-series fitted
%      rate constants themselves (s^-1), math-mode, e.g. "$1.2\times10^{-6}$"
% =============================================================
\newcommand{\ChTwoEa}{157.1}
\newcommand{\ChTwoEaErr}{5.3}
\newcommand{\ChTwoTc}{98.1}
\newcommand{\ChTwoTcErr}{8.6}
\newcommand{\ChTwoRtwo}{0.989}
\newcommand{\ChTwoN}{12}
\newcommand{\ChTwoPermEa}{10192}
\newcommand{\ChTwoPermEaErr}{1882}
\newcommand{\ChTwoPermRtwo}{0.936}
\newcommand{\ChTwoPermN}{10}
\newcommand{\ChTwoPermLamMin}{$2.7\times10^{-5}$}
\newcommand{\ChTwoPermLamMax}{$6.3\times10^{-5}$}

% Auto-generated by Chapter2/Figures/fig_winsor_mdd_free.py -- do not edit by hand.
\newcommand{\MDDFreeTableRows}{%
1 & 1 & 84.06 & -645.7 & 23.33 & 0.0 & 0.79 & 6.61 & 200 & 55.0 & 59.5 & 179 & 60--60 & 55.0--55.0 \\
2 & 3 & 19.76 & -834.2 & 3.24 & 2.1 & 1.03 & 1.33 & 200 & 97.8 & 65.0 & 178 & 63--67 & 87.5--100.0 \\
3 & 5 & 8.50 & -940.5 & 1.60 & 2.8 & 1.12 & 5.87 & 190 & 74.0 & 60.4 & 173 & 57--64 & 57.5--92.5 \\
4 & 7 & 8.05 & -939.4 & 1.58 & 2.8 & 1.12 & 5.59 & 190 & 74.4 & 60.4 & 173 & 57--64 & 60.0--95.0 \\
5 & 9 & 8.05 & -930.9 & 1.71 & 2.8 & 1.12 & 5.59 & 190 & 74.4 & 60.4 & 173 & -- & -- \\
6 & 11 & 8.05 & -922.5 & 1.85 & 2.8 & 1.12 & 5.59 & 190 & 74.4 & 60.4 & 173 & -- & -- \\
7 & 13 & 8.05 & -914.1 & 2.03 & 2.8 & 1.12 & 5.59 & 190 & 74.4 & 60.4 & 173 & -- & -- \\
8 & 15 & 8.05 & -905.6 & 2.24 & 2.8 & 1.12 & 5.59 & 190 & 74.4 & 60.4 & 173 & -- & -- \\
9 & 17 & 8.05 & -897.2 & 2.51 & 2.8 & 1.12 & 5.59 & 190 & 74.4 & 60.4 & 173 & -- & -- \\
10 & 19 & 8.05 & -888.7 & 2.84 & 2.8 & 1.12 & 5.59 & 190 & 74.4 & 60.4 & 173 & -- & -- \\
}
\newcommand{\MDDFreeBestBicN}{3}
\newcommand{\MDDFreeBestSpecN}{2}
\newcommand{\MDDFreeBestSpecChi}{1.33}
\newcommand{\MDDFreeTmaxMin}{59}
\newcommand{\MDDFreeTmaxMax}{65}
\newcommand{\MDDFreeTpeakMin}{55}
\newcommand{\MDDFreeTpeakMax}{98}
\newcommand{\MDDFreeNmax}{10}

\newcommand{\He}[1]{\textsuperscript{#1}He}

\title{Free multi-domain diffusion models and the thermal history of\\
Win-06-01b goethite from the water-saturated step heat HT-Win-06-01b}
\author{B. J. Thyer}
\date{\today}

\begin{document}
\maketitle

\begin{abstract}
We fit the $^3$He release of the water-saturated step-heating experiment on
proton-irradiated Win-06-01b goethite with multi-domain diffusion (MDD) models
of $N = 1$--\MDDFreeNmax{} spherical domains whose $^3$He fractions and sizes are all
free, following \citet{farleyflowers2012}. Each domain structure is then used to
fit a burial-and-exhumation thermal history to the bulk-normalized
$^4$He/$^3$He spectrum and to predict the (U--Th)/He age. The domain structure
stabilizes at $N=4$. About 97\% of the $^3$He resides in one domain slightly
larger than the Stage~1 reference domain ($r/r_0=1.12$). Three small domains
holding 2.8\% in total account for the high early diffusivities and the
gradual approach to the Stage~1 line. Larger $N$ only duplicate the dominant
domain. With this structure, burial paths peak at
58--64\,\textdegree C between $\sim$60 and 100~Ma, consistent with regional
apatite fission-track constraints. The He age requires crystallization at or
after $\sim$195~Ma.
\end{abstract}

% -------------------------------------------------------------
\section{Motivation and data}

Win-06-01b is dense colloform goethite collected as float in the Flinders
Ranges, South Australia, with a (U--Th)/He age of 130--180~Ma
\citep{waltenberg2012}. Apatite fission-track studies in the region record
cooling from roughly 50--70\,\textdegree C after the mid-Cretaceous
\citep{gibsonstuwe2000,mitchell1998}, and Triassic intramontane basins less
than 1.5~km thick indicate that parts of the ranges were near the surface in the
Early Mesozoic \citep{mitchell1998}. A proton-irradiated aliquot was step-heated
in a sealed silica ampoule under water-saturated conditions (HT-Win-06-01b; 34
steps, 170--250\,\textdegree C), releasing 14.8\% of its $^3$He before final
fusion. The Stage~1 regression of that experiment defines the reference
Arrhenius line used here
($E_a=\ChTwoEa$~kJ\,mol$^{-1}$, $\ln(D_0/a_0^2)=16.16$ with $D_0/a_0^2$ in s$^{-1}$).

Earlier versions of this analysis split the sample into a high-retentivity
domain pinned to the Stage~1 line and a low-retentivity assemblage holding at most 7.5\% of the $^3$He. Here that distinction is dropped: all
domains are equivalent, and each has a free $^3$He fraction and size.

% -------------------------------------------------------------
\section{Methods}

\subsection{Domain model}
The sample is represented by $N$ non-interacting spherical domains that share
the Stage~1 activation energy and frequency factor. Domain $j$ holds a
fraction $f_j$ of the uniformly distributed, proton-induced $^3$He
($\sum_j f_j=1$) and has radius $r_j$ relative to the reference radius $r_0$
implied by the Stage~1 regression, so that
\begin{equation}
  \left(\frac{D}{a^2}\right)_j = \left(\frac{r_j}{r_0}\right)^{-2}\frac{D_0}{a_0^2}\exp\!\left(-\frac{E_a}{RT}\right).
  \label{eq:domain}
\end{equation}
A model with $N$ domains has $k=2N-1$ free parameters. Sizes are bounded by
$10^{-6}\le r_j/r_0\le10^{3}$.

\subsection{Laboratory release of $^3$He}
For a heating schedule of steps $i$ with temperatures $T_i$ and durations
$\Delta t_i$, the reference domain accumulates the non-dimensional time
\begin{equation}
  \tau_{0,i} = \sum_{m\le i}\frac{D_0}{a_0^2}\exp\!\left(-\frac{E_a}{RT_m}\right)\Delta t_m ,
  \label{eq:tau0}
\end{equation}
and domain $j$ reaches $\tau_{0,i}(r_j/r_0)^{-2}$. An initially uniform sphere
releases the cumulative fraction \citep{crank1979}
\begin{equation}
  F(\tau) = 1-\frac{6}{\pi^2}\sum_{k=1}^{\infty}\frac{1}{k^2}\exp\!\left(-k^2\pi^2\tau\right),
  \label{eq:Fsphere}
\end{equation}
so the model cumulative and fractional $^3$He yields of the sample are
\begin{equation}
  \Sigma F_i = \sum_{j=1}^{N} f_j\,F\!\left(\tau_{0,i}\,(r_j/r_0)^{-2}\right),
  \qquad \Delta F_i = \Sigma F_i - \Sigma F_{i-1}.
  \label{eq:mixture}
\end{equation}

\subsection{Fitting the domain structure to $^3$He}
Following \citet[supplement~A2]{farleyflowers2012}, the fractional and cumulative
yields are weighted equally, and $\{f_j, r_j\}$ minimize
\begin{equation}
  S = \sum_{i=1}^{34}\left(\Delta F_i^{\mathrm{mod}}-\Delta F_i^{\mathrm{obs}}\right)^2
    + \sum_{i=1}^{34}\left(\Sigma F_i^{\mathrm{mod}}-\Sigma F_i^{\mathrm{obs}}\right)^2 ,
  \label{eq:S}
\end{equation}
with yields expressed as fractions of the total $^3$He, including the fusion
step. Models are built up one domain at a time: the $N$-domain search starts from
the $(N-1)$-domain solution plus one new domain trialled at 15 sizes and two
fractions, together with ten random starts. The fit therefore cannot get worse
as $N$ increases. To compare models of different complexity, we use the
Bayesian information criterion for Gaussian residuals of unknown variance,
\begin{equation}
  \mathrm{BIC} = n\ln(S/n) + k\ln n,\qquad n = 68 .
  \label{eq:bic}
\end{equation}
The Arrhenius and delta plots are not fitted. They follow from applying the
single-sphere inversion of \citet{fechtigkalbitzer1966} to the model
cumulative release. Their misfit is reported as a reduced $\chi^2$ against the
bounding-method uncertainties in $\ln(D/a^2)$.

\subsection{Radiogenic $^4$He in nature}
Each domain produces $^4$He uniformly at rate $P$ and loses it by diffusion
with a zero-concentration boundary,
\begin{equation}
  \frac{\partial C}{\partial t}=\frac{D}{a^2}\,\frac{1}{x^2}\frac{\partial}{\partial x}\!\left(x^2\frac{\partial C}{\partial x}\right)+P,
  \qquad C(1,t)=0,\quad C(x,t_\mathrm{form})=0 ,
  \label{eq:pde}
\end{equation}
where $x$ is radial position normalized to the domain radius and
$t_\mathrm{form}$ is the crystallization age. With
$C=\sum_k A_k(t)\sin(k\pi x)/x$ \citep{carslawjaeger1959}, each mode evolves as
\begin{equation}
  \frac{\mathrm{d}A_k}{\mathrm{d}t}=-k^2\pi^2\frac{D}{a^2}A_k+c_kP,\qquad c_k=\frac{2(-1)^{k+1}}{k\pi},
  \label{eq:modes}
\end{equation}
which over an interval $\Delta t$ (2~Myr) at constant temperature has the exact solution
\begin{equation}
  A_k \leftarrow A_k e^{-\lambda_k}+c_kP\,\Delta t\,\frac{1-e^{-\lambda_k}}{\lambda_k},
  \qquad \lambda_k=k^2\pi^2\frac{D}{a^2}\Delta t .
  \label{eq:update}
\end{equation}
We use 8000 modes. The $^4$He content per unit volume of domain $j$ is
$N_j=3\sum_kA_{k,j}(-1)^{k+1}/(k\pi)$. Taking U and Th to be distributed like
the $^3$He and neglecting parent decay ($\sim$3\% over 200~Myr), the bulk He
age is
\begin{equation}
  t_\mathrm{He}=\frac{1}{P}\sum_{j}f_j\,N_j .
  \label{eq:age}
\end{equation}

\subsection{Laboratory degassing of the natural profiles}
In the laboratory each mode decays as $\exp(-k^2\pi^2\tau)$. The reference
times $\tau_{0,i}$ used for the $^4$He prediction are solved from the
\emph{measured} cumulative $^3$He by inverting Eq.~\ref{eq:mixture}. The
predicted spectrum is therefore independent of the heating schedule and of
step-to-step kinetic scatter:
\begin{equation}
  \frac{R_{\mathrm{step},i}}{R_\mathrm{bulk}}=\frac{\Delta N_{4,i}/N_{4,\mathrm{tot}}}{\Delta F_i^{\mathrm{obs}}},
  \qquad
  \Delta N_{4,i}=\sum_jf_j\left[N_j(\tau_{0,i-1})-N_j(\tau_{0,i})\right] ,
  \label{eq:ratio}
\end{equation}
where $N_j(\tau)$ is the $^4$He remaining in domain $j$ after laboratory time
$\tau(r_j/r_0)^{-2}$ and $N_{4,\mathrm{tot}}=\sum_j f_jN_j(0)$ includes the
gas released at fusion, consistent with the measured bulk ratio.

\subsection{Thermal histories}
The thermal history is piecewise linear through $(t_\mathrm{form},T_s)$,
$(t_\mathrm{peak},T_\mathrm{max})$, $(50~\mathrm{Ma},T_s)$ and $(0,T_s)$ with
$T_s=20$\,\textdegree C. The goethite forms at the surface, is buried and heated
to a peak temperature, and is exhumed by the Paleogene, when Paleogene
sediments in Flinders Ranges valleys show that basement was exposed
\citep{mitchell1998}. For each domain structure and each
$t_\mathrm{form}\in\{180,190,200,210\}$~Ma, $(t_\mathrm{peak},T_\mathrm{max})$
are optimized by Nelder--Mead against
\begin{equation}
  \chi^2=\sum_i\left(\frac{R_i^{\mathrm{mod}}-R_i^{\mathrm{M}}}{\sigma_i}\right)^2,
  \qquad\sigma_i=\max\!\left[\tfrac12\left(R_i^{\mathrm{UB}}-R_i^{\mathrm{LB}}\right),\,0.02\right],
  \label{eq:chi2}
\end{equation}
over the 24 steps that have lag-correction bounds and more than 0.1~pccSTP of
$^3$He. M, LB and UB are the central, lower-bound and upper-bound
bulk-normalized ratios of the bounding method. We also map $\chi^2$ on a grid
of $t_\mathrm{peak}$ (2.5~Myr steps) and $T_\mathrm{max}$ (1\,\textdegree C
steps) for $t_\mathrm{form}=180$, 195 and 210~Ma. The joint 1$\sigma$ region is
$\chi^2\le\chi^2_{\min}+3.53\max(1,\chi^2_\nu)$, and the predicted He age is
contoured at the limits of the measured range. For each crystallization age, we draw 300 burial
paths from the gridded posterior $p\propto\exp[-(\chi^2-\chi^2_{\min})/(2s)]$,
with $s=\max(1,\chi^2_\nu)$, jittering each draw within its grid cell. With two
free parameters on a dense grid, this is equivalent to Markov-chain Monte Carlo
sampling of the posterior.

% -------------------------------------------------------------
\section{Results}

\subsection{Domain structure from $^3$He}
Table~\ref{tab:summary} lists the fits for all $N$, and Fig.~\ref{fig:domains}
shows the fitted domain spectra.
\begin{itemize}
  \item \textbf{$N=1$.} A single domain cannot fit the release (Arrhenius $\chi^2_\nu=23$).
  \item \textbf{$N=2$.} The second domain captures the high early diffusivities but leaves a poor Arrhenius fit ($\chi^2_\nu=3.2$).
  \item \textbf{$N=3$.} The yield misfit falls by more than half and the Arrhenius $\chi^2_\nu$ falls to 1.6.
  \item \textbf{$N=4$.} The smallest domain splits in two, and the fit improves slightly further (Arrhenius $\chi^2_\nu=1.58$; BIC indistinguishable from $N=3$).
  \item \textbf{$N\ge5$.} Every added domain coincides in size with the dominant domain. The fit is unchanged and BIC rises.
\end{itemize}
We therefore adopt $N=4$ as the complete description of the $^3$He data. Its
domains are:
\begin{itemize}
  \item 0.54\% of the $^3$He at $r/r_0=6.6\times10^{-4}$;
  \item 0.66\% at $2.6\times10^{-3}$;
  \item 1.6\% at 0.033;
  \item 97.2\% at 1.12.
\end{itemize}
The three small domains are fully degassed during the experiment. The dominant
domain releases 12\% of its $^3$He, at a diffusivity 0.80 times the Stage~1
line. The Stage~1 line is therefore a slightly fast description of the dominant
domain: it was regressed on bulk-inverted diffusivities that still carry a
small contribution from the fast domains. With the dominant domain free to
differ from the Stage~1 reference, the model follows the delta plot closely,
including the gradual decline between $\Sigma F_{^3\mathrm{He}}=0.04$ and 0.07
(Fig.~\ref{fig:delta}). Only the prograde 240 and 250\,\textdegree C steps in
that interval lie above the model, by $\sim$1--1.7$\sigma$.
Figure~\ref{fig:supp} shows the fitted yields and the resulting Arrhenius and
delta plots for $N=1$--5.

\begin{table}[htbp]
\centering
\caption{Summary of the free multi-domain models. RMS is the yield misfit
($\times10^{-4}$, fraction of total $^3$He); ``fast'' is the $^3$He fraction in
domains with $r/r_0<0.1$; $r_\mathrm{main}$ is the size of the dominant domain.
Thermal columns give the best burial path with a predicted He age of
130--180~Ma; the last two columns give the ranges of $T_\mathrm{max}$ (\textdegree C)
and $t_\mathrm{peak}$ (Ma) within the joint 1$\sigma$ region of the $\chi^2$
maps that also satisfy the He age (maps computed for $N\le4$).}
\label{tab:summary}
\resizebox{\textwidth}{!}{%
\begin{tabular}{rrrrrrrrrrrrll}
\toprule
$N$ & $k$ & RMS & BIC & Arrh.\ $\chi^2_\nu$ & fast (\%) & $r_\mathrm{main}/r_0$ &
$^4$He/$^3$He $\chi^2_\nu$ & $t_\mathrm{form}$ & $t_\mathrm{peak}$ & $T_\mathrm{max}$ &
$t_\mathrm{He}$ & $T_\mathrm{max}$ (1$\sigma$) & $t_\mathrm{peak}$ (1$\sigma$)\\
\midrule
\MDDFreeTableRows
\bottomrule
\end{tabular}}
\end{table}

\begin{figure}[htbp]
\centering
\includegraphics[width=0.72\textwidth]{mdd_free_domains.pdf}
\caption{Domain spectra of the best-fitting models. Symbol area scales with
the square root of the $^3$He fraction; the dashed line marks the Stage~1
reference domain. The three small domains are unchanged from $N=4$ onward;
domains added beyond $N=4$ coincide in size with the dominant domain.}
\label{fig:domains}
\end{figure}

\begin{figure}[htbp]
\centering
\includegraphics[width=0.72\textwidth]{ch2_mdd_delta_N4.pdf}
\caption{Delta plot of the measured $^3$He diffusivities (departure from the
Stage~1 line) and the curve implied by the four-domain fit to the $^3$He yields
(the delta plot itself is not fitted). The shaded band marks
$\Sigma F_{^3\mathrm{He}}=0.04$--0.07, where the prograde 220--250\,\textdegree C
steps show elevated diffusivities.}
\label{fig:delta}
\end{figure}

\subsection{Thermal history from $^4$He/$^3$He}
With the four-domain structure fixed, burial histories peaking at
$\approx$60\,\textdegree C fit the $^4$He/$^3$He spectrum best
(Figs.~\ref{fig:spectrum} and \ref{fig:mappaths}). The spectrum constrains the
peak temperature tightly, while peak time trades off against crystallization
age (Table~\ref{tab:posterior}):
\begin{itemize}
  \item \textbf{180~Ma.} Posterior paths peak at 59--63\,\textdegree C between 77.5 and 60~Ma.
  \item \textbf{195~Ma.} Paths peak at 59--64\,\textdegree C between 92.5 and 72.5~Ma.
  \item \textbf{210~Ma.} Paths peak at 58--64\,\textdegree C between 102.5 and 80~Ma.
\end{itemize}
The He-age isochrons then select among these. All posterior paths for
crystallization at 180~Ma and 97\% of those at 195~Ma predict He ages within the
measured 130--180~Ma range. None do for crystallization at 210~Ma, where every path
predicts an age older than 180~Ma. The surviving paths peak in the Late
Cretaceous, within or just after the apatite fission-track window. Their peak
temperatures fall within the 50--70\,\textdegree C range of
\citet{gibsonstuwe2000} and at the $<$60\,\textdegree C limit reported near
Copley \citep{mitchell1998}.

The absolute quality of the spectrum fit is modest ($\chi^2_\nu\approx5.2$--5.8).
The model underpredicts the ratios between $\Sigma F_{^3\mathrm{He}}=0.04$ and
0.07 and overpredicts the long 250\,\textdegree C step. A two-domain structure
fits the spectrum considerably better ($\chi^2_\nu\approx1.3$; peak
$\approx$65\,\textdegree C at 80--100~Ma). However, it reproduces the $^3$He
Arrhenius array poorly (Fig.~\ref{fig:supp}), so we do not adopt it.

\begin{table}[htbp]
\centering
\caption{Posterior burial paths for the four-domain structure (300 samples per
crystallization age; 16th--84th percentiles).}
\label{tab:posterior}
\begin{tabular}{rrrrr}
\toprule
$t_\mathrm{form}$ (Ma) & best $\chi^2_\nu$ & $T_\mathrm{max}$ (\textdegree C) & $t_\mathrm{peak}$ (Ma) & He age 130--180~Ma\\
\midrule
180 & 5.81 & 59--63 & 60.0--77.5 & 100\% \\
195 & 5.50 & 59--64 & 72.5--92.5 & 97\% \\
210 & 5.20 & 58--64 & 80.0--102.5 & 0\% \\
\bottomrule
\end{tabular}
\end{table}

\begin{figure}[htbp]
\centering
\includegraphics[width=0.72\textwidth]{ch2_mdd_spectrum_N4.pdf}
\caption{Bulk-normalized $^4$He/$^3$He spectrum (black, bounding-method LB--UB
boxes) and the prediction of the best age-consistent burial path for the
four-domain structure.}
\label{fig:spectrum}
\end{figure}

\begin{figure}[htbp]
\centering
\includegraphics[width=\textwidth]{ch2_mdd_map_paths_N4.pdf}
\caption{Thermal-history constraints for the four-domain structure, with one
column per assumed crystallization age (180, 195 and 210~Ma). Top row: reduced
$\chi^2$ of the $^4$He/$^3$He spectrum as a function of the time and magnitude
of peak temperature. The colour scale is anchored to this model's minimum. White
outlines show the joint 1$\sigma$ region, dashed lines the predicted bulk He age
at 130 and 180~Ma, diamonds the best fit, and orange boxes the apatite
fission-track peak window \citep{gibsonstuwe2000,mitchell1998}. Bottom row:
300 burial paths drawn from the gridded posterior for each crystallization age.
Blue paths predict He ages within 130--180~Ma; grey paths do not. The best fit
is shown in black.}
\label{fig:mappaths}
\end{figure}

\begin{landscape}
\begin{figure}[htbp]
\centering
\includegraphics[width=\linewidth]{mdd_free_N4.pdf}
\caption{Diagnostic summary of the four-domain model: (a) fitted fractional
and (b) cumulative $^3$He yields; (c) Arrhenius and (d) delta plots implied by
the fit; (e) domain spectrum; (f) $^4$He/$^3$He spectrum with the best
age-consistent burial path; (g) best path for each crystallization age; (h)
parameter summary, including the fraction of each domain degassed during
the experiment.}
\label{fig:diagN4}
\end{figure}
\end{landscape}

% -------------------------------------------------------------
\section{Discussion}

\paragraph{Robust conclusions.} Across all domain structures, the dominant
domain holds $\approx$97\% of the $^3$He and degasses slightly more slowly
than the Stage~1 line. Fast domains hold only 2--3\%. With the
four-domain structure, burial paths peak at 58--64\,\textdegree C in the Late
Cretaceous, followed by exhumation, and the He age requires crystallization at
or after $\sim$195~Ma. This thermal history is consistent with the apatite fission-track
record of 50--70\,\textdegree C before cooling after $\sim$85~Ma
\citep{gibsonstuwe2000} and with the shallow Triassic basins and low-rank coals
of the Flinders Ranges \citep{mitchell1998}.

\paragraph{Remaining misfit.} The four-domain structure leaves systematic
misfit in the $^4$He/$^3$He spectrum. Its intermediate domain ($r/r_0\approx0.03$,
1.6\% of the $^3$He) degasses at $\Sigma F_{^3\mathrm{He}}\approx0.02$--0.06. With
the common activation energy, it would have lost most of its radiogenic $^4$He in
nature. That depresses the predicted ratios where the measured ratios are
highest. Diffusion-only remedies that preserve the $^3$He fit are:
\begin{itemize}
  \item a higher activation energy for the intermediate domain, which the retrograde segments of the experiment could test;
  \item domain-dependent $^4$He production, for example from U--Th zoning across the colloform bands.
\end{itemize}
Each adds one parameter. A finite non-diffusive release during the first heating
to 220--250\,\textdegree C remains possible but is no longer needed to explain
the $^3$He data.

\paragraph{Caveats.}
\begin{itemize}
  \item The temperature scale rests on extrapolating the Stage~1 activation energy from 200--250\,\textdegree C to near-surface temperatures ($\pm\sim$4\,\textdegree C near 60\,\textdegree C, not propagated here).
  \item The exhumation time (50~Ma) and surface temperature (20\,\textdegree C) are fixed.
  \item Domains that barely degas in the experiment are constrained only through their summed fraction.
\end{itemize}

\bibliographystyle{plainnat}
\bibliography{../../references}

% -------------------------------------------------------------
\appendix
\section{Supplementary figure}
Diagnostic figures and $\chi^2$ maps for the other domain structures are written
to \texttt{Figures/generated/mdd\_free/absolute/}.

\begin{figure}[p]
\centering
\includegraphics[width=\textwidth,height=0.88\textheight,keepaspectratio]{ch2_mdd_3he_fits_N1-5.pdf}
\caption{Fits to the $^3$He data for $N=1$--5 domains (rows). Columns: fitted
fractional and cumulative yields, and the Arrhenius and delta plots implied by
each fit. Hollow symbols are low-signal steps or steps without lag bounds.}
\label{fig:supp}
\end{figure}

\end{document}
